\documentclass[11pt,a4paper]{article}
% Shared preamble for RuPhO English problem statements (match T1 2026 style).
% Use: \input{latex/rupho_en_preamble.tex} after \documentclass.
\usepackage[margin=1in]{geometry}
\usepackage{amsmath,amssymb}
\usepackage{graphicx}
\usepackage{float}
\usepackage{enumitem}
\usepackage{microtype}
\usepackage{caption}
\usepackage{tabularx}
\usepackage{hyperref}

\captionsetup{font=small,labelfont=bf,skip=6pt}

\setlist[itemize]{leftmargin=1.2cm}
\setlist[enumerate]{leftmargin=1.2cm}

% One outer box per subproblem: [id | statement | points] (no inner rules)
\newcommand{\problembox}[3]{%
  \par\addvspace{0.42em}%
  \noindent
  \setlength{\fboxsep}{7.5pt}%
  \setlength{\fboxrule}{0.95pt}%
  \fbox{%
    \begin{minipage}{\dimexpr\linewidth-2\fboxsep-2\fboxrule\relax}
    \begin{tabularx}{\linewidth}{@{}>{\raggedright\arraybackslash\bfseries}p{2.1em} @{\hspace{0.9em}} >{\raggedright\arraybackslash}X @{\hspace{0.9em}} >{\raggedleft\arraybackslash\bfseries}p{2.4em}@{}}
    #1 & #2 & #3
    \end{tabularx}
    \end{minipage}%
  }%
  \par\addvspace{0.32em}%
}

\title{\textbf{Propagation of detonation waves}}
\author{\normalsize Y23 (2023) --- English translation}
\date{}
\begin{document}
\maketitle

Combustion is an exothermic chemical reaction (fuel + oxidizer), accompanied by positive specific heat release. Under certain conditions, it can be self-sustaining, in contrast to endothermic (accompanied by the absorption of heat) reactions. For example, when a source of ignition occurs in a pre-mixed gas mixture of fuel and oxidizer, a combustion wave can form. Detonation is a combustion mode in which a shock wave propagates through a substance, initiating chemical combustion reactions that support the movement of the shock wave.

A flame front is a narrow reaction zone of a propagating flame in which combustion occurs. This task asks you to study the propagation of detonation shock waves in a gas.

Let the gas move through a thermally insulated straight pipe of constant cross-section. The gas can be divided into two areas. In one region, the gas is motionless, its pressure is $p_0$, density is $\rho_0$. In another region, all elements of the gas move at the same speed, gas pressure $p_1$, density $\rho_1$. The speed of movement of the boundary between these regions $c$ will be called the shock wave speed. The acceleration of free fall, as well as the thermal conductivity of the gas and its friction against the walls of the vessel can be neglected. Note: solving the problem will be significantly simplified if we reduce the gas movement to stationary. During stationary motion, the gas parameters at each point in space remain constant in time.

\section*{Part A. Shock wave in gas (3.0 points)}

A wave in which sharp jumps in pressure and density occur near the wave front is called a shock wave. Let us consider a one-dimensional flow of an ideal gas with the adiabatic exponent $\gamma$. In this part of the problem, you are asked to analyze the relationships between the parameters $p_0$, $\rho_0$, $p_1$ and $\rho_1$ in the shock wave.

\begin{figure}[H]
    \centering
    \includegraphics[width=0.72\textwidth]{T4_EN_assets/fig1.jpg}
    \caption{Figure 1.}
\end{figure}

\problembox{A1}{Write down a system of laws of conservation of mass and momentum for the volume of gas passing through the wave front of the shock wave. Determine from it the speed of movement of the wave front $c$.}{1.00}

Since there are no dissipations in the system, the energy in it is also conserved.

\problembox{A2}{Write down the law of conservation of energy for an element of the volume of gas passing through the wave front. Obtain from it an additional equation relating the parameters $p_0$, $\rho_0$, $p_1$, $\rho_1$, the speed of movement of the wave front $c$ and the adiabatic exponent $\gamma$.}{1.00}

\problembox{A3}{Combining the results obtained in steps A1 and A2, express the pressure ratio $p_1/p_0$ in terms of the adiabatic exponent $\gamma$ and the compaction coefficient $k=\rho_1/\rho_0$. This relationship is called the adiabatic equation for a shock wave.}{0.50}

\problembox{A4}{Determine the maximum possible value of the compaction coefficient $k_{max}$. Express your answer in terms of $\gamma$. Calculate the answer for an ideal diatomic gas.}{0.50}

\section*{Part B. Propagation of detonation combustion wave (4.0 points)}

When heated, some gases begin a chemical reaction and ignite. A shock wave can then propagate through the gas, and combustion occurs in a very thin layer near the wave front. In this part of the problem we will consider a detonation wave. The main difference between this part of the problem and part A is that when the mass $\Delta{m}_0$ enters the motion of the system (when this mass passes through the wave front), a chemical reaction occurs and heat $\Delta{Q}=q\Delta{m}_0$ is released, where $q$ is a known quantity called the specific thermal effect of the chemical reaction.

\begin{figure}[H]
    \centering
    \includegraphics[width=0.72\textwidth]{T4_EN_assets/fig2.jpg}
    \caption{Figure 2.}
\end{figure}

Since during a chemical reaction the composition of the substance changes, the adiabatic index $\gamma$ also changes. Let us denote by $\gamma_0$ and $\gamma_1$ the adiabatic exponents of the stationary gas and the gas after the passage of the wave front, respectively.

\problembox{B1}{Write down the law of conservation of energy for a detonation wave. Obtain from it an equation relating the parameters $p_0$, $\rho_0$, $p_1$, $\rho_1$, the speed of movement of the wave front $c$, the specific thermal effect of the reaction $q$, as well as the adiabatic exponents $\gamma_0$ and $\gamma_1$.}{0.60}

Let us introduce the value of the specific volume $v=1/\rho$ for each state: $$v_0=\cfrac{1}{\rho_0},\qquad v_1=\cfrac{1}{\rho_1}. $$  $p_1(\rho_1)$,  for       $c$      s     $p_1(\rho_1)_\text{imp }$.   $p_1(\rho_1)$,      ,  and    $p_1(\rho_1)_\text{child}$.   in    , t. ,  : $$p_0\ll{p_1},\qquad \cfrac{p_0}{\rho_0}\ll{q}. $$

\problembox{B2}{The explicit dependence $p_1(v_1)_\text{}$ for a strong detonation wave can be represented as: $$p_1=\cfrac{q}{v_0}\cdot A $$ $A$.    $\gamma_1$ and  $v_1/v_0$. This relationship is called the detonation adiabatic equation.}{0.60}

\problembox{B3}{On one graph, qualitatively depict the dependence $p_1(v_1)_\text{}$, as well as the dependence $p_1(v_1)_\text{}$ for different values of the speed of movement of the wave front $c$.}{0.60}

When studying detonation processes, Chapman and Jouguet hypothesized that the detonation combustion wave propagates at the minimum possible speed at which the equation $p_1(v_1)_\text{}=p_1(v_1)_\text{}$ has a unique solution, called the Chapman-Jouguet point. This statement was later proven by Zeldovich and Doering.

\problembox{B4}{Determine the parameters $p_1$ and $v_1$ of the Chapman-Jouguet point. Express your answer in terms of $v_0$, $\gamma_1$, $q$.}{1.40}

\problembox{B5}{Determine the speed of movement of the wave front $c$. Express your answer in terms of $q$ and $\gamma_1$.}{0.20}

\problembox{B6}{Calculate the speed of movement of the wave front during a detonation wave in a mixture of oxygen and hydrogen if, when it is heated to high temperatures, the following chemical reaction occurs: $$ 2\mathrm{H}_2+ \mathrm{O} _2 = \mathrm{H}_2 \mathrm{O} + \mathrm{H}^{+} + \mathrm{HO}^{-}. $$       , in     and ,  in       $\mathrm{H}^{+}$ and $\mathrm{HO}^{-}$ (     ). V cm on 2 mol   mol  in  s  .    ,  $q \approx 15~\text{MJ}/\text{kg}$ for every kilogram of hydrogen consumed.}{0.60}

\section*{Part C. Luminous detonation (3.0 points)}

Under the influence of a strong electric field, the air is ionized and turns into plasma. A sufficiently strong field can be obtained by focusing laser radiation. The plasma, in turn, contains a large number of free charges and absorbs electromagnetic waves well, which leads to further heating of the plasma and propagation of the shock wave. On one side of the wave front there is plasma, on the other there is air. At the interface between air and plasma, laser radiation is absorbed and its energy is transferred to the plasma. Excessive pressure of heated air creates a shock wave propagating towards the laser radiation. This phenomenon is called light denotation wave. Ramsden and Savich proposed to describe such a wave as a detonation wave, in which the analogue of the specific thermal effect of the reaction $q$ is the value $$ q_{\text{}}  =\frac{I_0}{\rho_0 c}, $$where $I_0$ —  . V    , in    ,    $P = 30 ~\text{MW}$,   s    in   $r \approx 10^{-2}~ \text{cm}$,   $$ I \approx \frac{P}{\pi r^2}. $$         $E_\text{probe}=30~\text{kV}/\text{cm}$, causing air breakdown, so it is enough to form a light-detonation wave.

\begin{figure}[H]
    \centering
    \includegraphics[width=0.72\textwidth]{T4_EN_assets/fig3.jpg}
    \caption{Figure 3.}
\end{figure}

Next, take the following data: $p_0=1.0 \cdot10^5~\text{Pa}$ and $\rho_0=1.3~\text{kg}/\text{m}^3$ are the pressure and density of air at atmospheric pressure, respectively. The plasma equation of state can be described by the equation of state of an ideal gas with the adiabatic exponent $\gamma_1=4/3$. Work in the approximation of a strong wave, i.e. consider $p_0\ll{p_1}$ and $p_0c\ll{I_0}$.

Next, accept the following data:

\problembox{C1}{Obtain an equation for the light detonation adiabatic similar to point B2, that is, find the pressure $p_1$ in the plasma after the passage of the wave front. Express your answer in terms of $v_0$, $v_1$, $\gamma_1$, $I_0$.}{1.20}

\problembox{C2}{Determine the parameters ($p_1$ and $v_1$) of the Champin-Jouguet point for the light-detonation adiabatic. Express your answer in terms of $v_0$, $\gamma_1$, $I_0$.}{1.40}

\problembox{C3}{Determine the speed of movement of the wave front $c$ of the light detonation wave. Express your answer in terms of $\gamma_1$, $I_0$, $\rho_0$. Calculate its numerical value.}{0.40}

\vspace{1em}
\noindent\textit{Source:} \url{https://pho.rs/p/3487}

\end{document}